\paragraph{Qué términos deja el cambio de soporte.}
Descompongamos \(S_0\) en bloques \(B_0,C_0,C_0^{\mathsf T},D_0\) según \(I,\mathcal J\).
Sea \(U_0=\sum_{i\in F}p_i a_{i,I}a_{i,I}^{\mathsf T}\),
\(\mathcal M=B_0-U_0\), \(d_0=C_0\mathbf1\), \(\beta=Ra_c\),
\(z=\beta_I\), \(t=\beta_{\mathcal J}-\mathbf1\) y \(r=d_0+\mathcal M z=-C_0t-U_0z\).
El mínimo de la energía cuadrática del hijo es \(Q_*=d_0^{\mathsf T}\mathcal M^+d_0\).
Al insertar \(d_0=r-\mathcal Mz\) y usar \(S_0\beta=a_c\) se obtiene
\[
 Q_*=a_c^{\mathsf T}Ra_c-4+12p_c+E_Q,\qquad
 E_Q=z^{\mathsf T}U_0z-t^{\mathsf T}D_0t+r^{\mathsf T}\mathcal M^+r .
\]
En efecto \(D_0=(1-p_c)I+p_c\mathbf1\mathbf1^{\mathsf T}\) y
\(\mathbf1^{\mathsf T}D_0\mathbf1=4+12p_c\). Se tiene \(r=O_b(m^{-2})\) por entrada
y \(r^{\mathsf T}\mathcal M^+r=O_b(m^{-3})\) por la brecha y el número \(O_b(m)\) de líneas.

Sea \(\gamma_f=h_f^{(c)}-h_f\) el cambio de vacancia en una diagonal retenida.
Taylor de la entropía del apéndice~\ref{app:predictor}, con los pesos del hijo, da
\begin{align*}
 \Delta H-\log p_c&=\tfrac32\sum_{F\setminus\{c\}}p_i^2-\tfrac12 Q_*+A_3+O_b(m^{-3}),\\
 A_3&=\sum_{F\setminus\{c\}}p_i^3+\tfrac16\sum_{E'}p_i(y_i/p_i)^3
 +\tfrac32\sum_{E'}\frac{y_i^2}{1-p_i}+\tfrac16\sum_{f\in I_D}\gamma_f^3/h_f^2 .
\end{align*}
Aquí \(I_D\) recorre ambas familias diagonales retenidas. Los términos cuárticos cuestan
\(O_b(m^{-3})\) por \(\sum p_i=m\), \(p_i=O_b(m^{-1})\) y
\(|y_i/p_i|+|\gamma_f/h_f|=O_b(m^{-1})\). La ortogonalidad en el mínimo deja un error
cuadrático en las correcciones relativas \(O_b(m^{-2})\) a la respuesta lineal.

Para el cofactor escribamos \(\Pi_m=vv^{\mathsf T}/(2m)\), donde \(v\) vale \(1,-1,0,0\)
en las cuatro familias. Sean \(H_*=(R+\Pi_m)_{II}\), \(C_{\mathcal J}=D_0+(\Pi_m)_{\mathcal J\mathcal J}\),
\(V_{\mathcal J}=C_0+(\Pi_m)_{I\mathcal J}\) y \(\Pi'=v_Iv_I^{\mathsf T}/[2(m-1)]\).
En \(I\), \(S_*\) tiene entrada \(-p_i\) en cada par ordenado de líneas distintas
de \(i\in F\setminus\{c\}\). En forma análoga \(Y\) tiene entrada \(y_i\) en los pares de
\(i\in E'\). Sus restantes entradas son cero. El complemento de Schur da
\begin{align*}
 D_*&=S_*+Y+\Pi'-(\Pi_m)_{II}+V_{\mathcal J} C_{\mathcal J}^{-1}V_{\mathcal J}^{\mathsf T},\\
 Q_c&=\tfrac{m^2}4\{\operatorname{tr}[(H_*D_*)^2]-\operatorname{tr}(S_*^2)\},\\
 E_R&=\operatorname{tr}[(C_{\mathcal J}-I)(H_{\mathcal J}-I)],\qquad H_{\mathcal J}=(R+\Pi_m)_{\mathcal
 J\mathcal J}.
\end{align*}
Para la normalización del calibre sean \(S'_0\) la matriz del hijo y
\(\Delta_0=S'_0-(S_0)_{II}\). Su contribución se obtiene de
\[
 E_g=\log[m/(m-1)]-\log\det C_{\mathcal J}
   +\operatorname{tr}[(\Pi_m)_{II}\Delta_0]
   +\operatorname{tr}[H_*(\Pi'-(\Pi_m)_{II})].
\]
Como \(v_{\mathcal J}\perp\mathbf1\), \(\det C_{\mathcal J}=(1+3p_c)(1-p_c)^2(1-p_c+1/m)\).
Los logaritmos dan \(m^{-2}-p_c/m+6p_c^2+O_b(m^{-3})\). Las dos trazas valen
\(-(1-p_c)/m\) y \((m-1)/m^2+v_{\mathcal J}^{\mathsf T}R_{\mathcal J\mathcal J}v_{\mathcal J}/[2m(m-1)]\),
 respectivamente,
por \(S_0v=Rv=0\). Como \(R_{\mathcal J\mathcal J}=I+O_b(m^{-1})\), queda \(E_g=m^{-2}+6p_c^2+O_b(m^{-3})\).
El cubo del logaritmo del determinante deja \(C^{(3)}=-\operatorname{tr}(S_*^3)/6\).
En efecto, las filas de \(S_*\) tienen suma absoluta \(O_b(m^{-1})\);
las entradas de \(D_*-S_*\) son \(O_b(m^{-2})\) y las de \(H_*-I\) son \(O_b(m^{-1})\).
En la traza cúbica, un factor de \(H_*D_*-S_*\) cuesta así \(O_b(m^{-3})\) en total.
Las potencias de orden cuatro o superior tienen la misma cota por la norma de filas.
La traza pura \(\operatorname{tr}(S_*^2)/4\) aporta
\(3\sum_{F\setminus\{c\}}p_i^2/2\), pues hay seis pares ordenados en cada casilla retirada.

En la parte lineal de Schur, sea \(f_i^I\) la restricción de \(f_i\) a pares retenidos.
Las cuatro entradas diagonales de \(S_0R=I-\Pi_m\) sobre \(\mathcal J\) dan
\[
 \sum_{i\in F}p_i(f_i-f_i^I)
 =-\operatorname{tr}[(R-I)_{\mathcal J\mathcal J}]-1/m-p_cf_c .
\]
Al juntar la traza lineal, \(-Q_*/2\), las dos copias de
\(3\sum_{F\setminus\{c\}}p_i^2/2\) y \(-3\Delta\sum p_i^2\), los términos con \(w_i\)
son \(\sum_F p_iw_i-w_c-\sum_{E'}y_iw_i\) y queda \(p_cf_c-3p_c^2\).
Restar \(\delta_c^0\) convierte las dos sumas en las de la última línea de la fórmula
\begin{align*}
 L(B+c)-L(B)-\log p_c-\delta_c^0
   &=\Delta_c V+P_c+O_b(m^{-3}),\\
 P_c&=-E_g/2+Q_c/m^2+C^{(3)}+A_3-3\sum_{E'}y_i^2\\
 &\quad-3p_c^2+p_cf_c-(E_Q+E_R)/2\\
 &\quad+\sum_F p_i(1-t_{ic})w_i-\sum_{E'}p_iq_iw_i .
\end{align*}
